\section*{填写说明}

本节用于撰写前检查，正式提交时可根据比赛要求删除或保留。本文档按一份主技术文档组织，目录采用“项目概述--总体设计--核心机制--实现测试--创新部署--总结展望”的竞赛作品报告结构，不再拆分多个独立说明书。

\begin{enumerate}
  \item 文档应围绕 OralSage 参赛作品本身展开，重点说明系统如何回应 A3 赛题要求。
  \item 若技术文档用于网评或匿名评审，请勿在正文、页眉页脚、图片、接口地址、代码仓库链接和截图中暴露学校、学院、指导教师、队员姓名等身份信息。
  \item 文中所有 \softtodo{占位说明} 均需替换为项目实际内容；不适用的小节可以删除或合并。
  \item 图、表、接口、测试用例应有编号和文字说明，关键结论应尽量配套截图、日志、测试记录或真实生成样例。
  \item 第三方库、模型、数据集、开源代码、科大讯飞相关工具和 AI Coding 工具的使用，应在“系统实现”章节中说明来源、版本、许可和使用边界。
\end{enumerate}

\begin{table}[H]
  \centering
  \caption{本文档整体撰写大纲}
  \label{tab:document-outline}
  \begin{tabularx}{\linewidth}{p{0.22\linewidth}X}
    \toprule
    章节 & 撰写重点 \\
    \midrule
    项目概述 & 项目背景、赛题目标拆解、口腔医学学习场景、痛点分析、作品定位、总体解决方案和核心功能。 \\
    系统总体设计 & 系统架构，前端、后端、知识库和模型服务关系，以及“画像--规划--资源生成--学习--测评--反馈”的主业务流程。 \\
    多智能体协同机制 & Orchestrator、Planner、Resource Factory、Tutor、QA、Assessment、Memory/User Model 的职责和协作关系。 \\
    个性化学习画像与学习路径规划 & 初始测评、动态画像更新、学习目标/节奏/偏好/薄弱点建模、个性化学习路径和学习计划推荐。 \\
    个性化资源生成设计 & 讲义、讲稿、PPT/slides、思维导图、题库练习、阶段卷、TTS 音频，以及“讲义作为资源母本”的设计思路。 \\
    知识库构建与防幻觉机制 & 数据输入与处理、口腔医学 KAG 知识库架构、知识库构建方法、证据引用、内容校验、结构化输出、fallback 和安全过滤。 \\
    系统实现 & 技术栈、后端接口、前端页面、数据存储与文件组织、部署运行方式、第三方组件与工具说明。 \\
    测试与评估 & 测试目标与环境、功能测试、智能体链路测试、资源生成质量测试、学习测评测试、防幻觉测试、响应时间与稳定性测试。 \\
    作品创新点与优势分析 & 技术创新、工程创新、应用创新、实用价值、痛点解决和与传统方案对比。 \\
    部署运行与演示说明 & 运行环境、配置项、部署步骤和演示流程。 \\
    总结与展望 & 作品完成情况、当前不足和后续规划。 \\
    \bottomrule
  \end{tabularx}
\end{table}

\begin{table}[H]
  \centering
  \caption{A3 赛题要求与章节对应关系}
  \label{tab:a3-map-new}
  \begin{tabularx}{\linewidth}{p{0.32\linewidth}X}
    \toprule
    A3 赛题要求 & 对应章节 \\
    \midrule
    对话式学习画像自主构建与动态更新 & 个性化学习画像与学习路径规划 \\
    多智能体协同架构设计 & 多智能体协同机制 \\
    至少 5 类个性化资源生成 & 个性化资源生成设计 \\
    个性化学习路径规划与资源推送 & 个性化学习画像与学习路径规划、系统总体设计 \\
    智能辅导与多模态答疑 & 多智能体协同机制、系统实现 \\
    学习效果评估与动态调整 & 多智能体协同机制、测试与评估 \\
    防幻觉与内容安全过滤 & 知识库构建与防幻觉机制 \\
    可运行源码、演示视频、PPT 和配套文档 & 系统实现、部署运行与演示说明、测试与评估 \\
    \bottomrule
  \end{tabularx}
\end{table}
